\documentclass[12pt]{article}
\usepackage[margin=1in]{geometry}
\usepackage{hyperref}
\pagestyle{empty}
\begin{document}

\begin{flushright}
Huang Zhongchang\\
Independent Researcher\\
Nanjing, China\\
\today
\end{flushright}

\vspace{0.3in}

\noindent The Editor\\
Physical Review D\\
American Physical Society

\vspace{0.2in}

\noindent \textbf{Re: Submission to Physical Review D}

\vspace{0.1in}

\noindent Dear Editor,

I am submitting ``One Definition, One Universe: Gravity and Quantum Decoherence from Causal Structure'' to Physical Review D.

The paper shows that Newtonian gravity and quantum decoherence emerge from a single mechanism---bounded information capacity of discrete cells---using a single definition: a causal event is an asymmetric influence producing a determinate change (from which the 1-bit quantization follows as a theorem). The free-capacity fraction $q$ obeys a unified telegraph equation. In the static limit $\nabla^2(\ln 1/q)=0$, the solution $q(r)=e^{-GM/rc^2}$ recovers the Newtonian potential, and the emergent metric $g_{00}=-q^2$, $g_{rr}=1/q^2$ yields PPN parameters $\beta=\gamma=1$, consistent with solar system tests. The $\gamma\to0$ vacuum limit dynamically selects Lorentz invariance. An H-theorem is proved via a Lyapunov functional. The decomposition $\nabla^2 q/q=|\nabla q|^2/q^2$ unifies diffusion (quantum decoherence) and self-organization (gravitational collapse) as two regimes of one process.

The quantum-mechanical foundation builds on Zilly (2026, arXiv:2603.11900). The classical gravitational derivation (Sections~3--7) depends only on Definition~C and structural choices, independently of the QM foundation.

Limitations are stated explicitly: only the Newtonian limit is recovered (not full GR), the volume-scaling $S\sim R^3$ is in tension with holographic area scaling, $\gamma(q)$ is a modeling ansatz, $d=3$ is empirical input, and matter coupling is proposed but not derived. A tiered-assumption classification (Table~I) is applied to both this framework and Bassani \& Magueijo (PRD 111, 103529, 2025), providing a transparent accounting of all axioms, empirical inputs, structural choices, and free parameters---a methodological contribution in its own right.

I suggest referees with expertise in emergent gravity, quantum foundations, and causal set theory. The manuscript is not under consideration elsewhere.

\vspace{0.15in}
\noindent Respectfully,

\vspace{0.3in}
\noindent Huang Zhongchang

\end{document}
